\documentclass[10pt]{beamer}
\usefonttheme{professionalfonts}
%\usetheme{Rochester}
%\usetheme{AnnArbor}
\usetheme{Szeged}
\useoutertheme{infolines}
\usecolortheme{dolphin}
\usepackage{graphicx}
% \usepackage[english]{babel}
\usepackage{hyperref}
\usepackage{multicol}
\usepackage{color}
%\usepackage{smartdiagram}

%\definecolor{carmine}{rgb}{0.59, 0.0, 0.09}
\usepackage{amsmath}

\usepackage{graphicx}
\usepackage{tikz}
\usetikzlibrary{arrows.meta}
\usetikzlibrary{shapes.callouts}
\usepackage{multirow}

\usepackage[utf8]{inputenc}

\usepackage{enumerate}
\usepackage[shortlabels]{enumitem}
\usepackage{amsmath,amsthm,amssymb,amsfonts,mathrsfs,latexsym,stmaryrd}
\usepackage{tcolorbox}

\usepackage{color}
\usepackage{amssymb, amsmath, amsbsy}
\usepackage{adjustbox,lipsum}
\usepackage{placeins}

\usepackage{textpos}
\usepackage{wrapfig}
\usepackage[spanish, es-tabla]{babel}

\usepackage{caption}
\setlength{\parindent}{15pt}
%\usepackage[hyperfootnotes=false]{hyperref}

\usepackage{url}
\usepackage{float}
\usepackage{tikz}
\usetikzlibrary{patterns}

\title[ICS-162]{Econometr\'ia ICS-294\\ \small{Clase 17: Endogenidad y Variables Instrumentales}}

\author{Jocelyn Tapia Stefanoni}
\institute[USM]
{Universidad T\'ecnica Federico Santa Mar\'ia\\ % Your institution for the title page
    \medskip
    \textit{} % Your email address
    }
\logo{\includegraphics[width=3cm]{certamenes/logo_usm.png}}
\date{}
\newcommand\Fontvi{\fontsize{9}{7.2}\selectfont}

\begin{document}

\begin{frame}
    \titlepage % Print the title page as the first slide
\end{frame}

\section{Introducci\'on}

\begin{frame}{Introducci\'on}
    \begin{itemize}
        \item En las clases pasadas vimos que cuando omitimos una variable $x_{2}$ que se correlaciona con $x_{1}$, se viola el supuesto $E\left[u \mid x_{1}\right]=0$ y nuestro estimador MCO estar\'a sesgado.
        \item Cuando esto sucede, decimos que $x_{1}$ es {\color{blue}end\'ogena (est\'a correlacionada con el error)}.
        \item {\color{blue}Una variable explicativa es endógena cuando su valor está parcialmente determinado dentro del sistema por relaciones que también afectan a la variable dependiente, lo que se manifiesta estadísticamente como correlación con el término de error. }
        \item \textbf{\textquestiondown Qu\'e podemos hacer}?
        \item M\'etodo 1: variables instrumentales (IV).
        \item  M\'etodo 2: mínimos cuadrados en 2 etapas (MC2E).
    \end{itemize}
\end{frame}

\section{Endogeneidad y IV}

\begin{frame}{Causas de la endogeneidad}
    Generalmente observamos problemas de endogeneidad en una de estas tres situaciones:
    \begin{enumerate}
        \item Variables omitidas:
        \begin{itemize}
            \item Son variables que est\'an correlacionadas con las $x$ (los regresores) pero no las podemos incluir por ser inobservables.
        \end{itemize}
        \item Simultaneidad: $x \rightarrow y$ y $y \rightarrow x$.
        \item Error de medida en las $x$.
    \end{enumerate}
\end{frame}

\begin{frame}{Endogeneidad}
    El problema de la endogeneidad:
    \begin{figure}
        \centering
        \includegraphics[width=0.3\linewidth]{complemetario/Imagenes/Grafico4_1.jpeg}
    \end{figure}


    
    $$
    E\left[\hat{\beta}_{1}^{M C O}\right]=\beta_{1}+\delta_{1} \beta_{2}
    $$

    Con  variable omitida.
\end{frame}


\begin{frame}{Soluci\'on 1: Variables instrumentales}
    Una variable instrumental v\'alida es una variable $\mathbf{Z}$ que se correlaciona con $X$, pero no con el error.
    \begin{figure}
        \centering
        \includegraphics[width=0.5\linewidth]{complemetario/Imagenes/Grafico4_2.jpeg}
    \end{figure}
    Formalmente, un instrumento es v\'alido si cumple con \textbf{dos} supuestos clave:
    \begin{enumerate}
        \item {\color{blue}\textbf{Relevancia:} }$\operatorname{Cov}(X, Z) \neq 0$. La variable $Z$ afecta a $X$. Como observamos $Z$ y $X$, este supuesto es testeable.
        \item  {\color{blue}\textbf{Exogeneidad o restricci\'on de exclusi\'on:}} $\operatorname{Cov}(Z, U)=0$. 
        \begin{itemize}
            \item El instrumento no se correlaciona con el error.
            \item $E[Y \mid X]=E[Y \mid X, Z]$: el instrumento $Z$ no afecta directamente a $Y$, s\'olo lo hace a trav\'es de su efecto en $X$.
        \end{itemize}
    \end{enumerate}
\end{frame}

\begin{frame}{Intuici\'on}
    \begin{itemize}
        \item La intuici\'on es que $X$ tiene dos componentes: una parte ex\'ogena y una parte end\'ogena (que est\'a correlacionada con $u$).
        \item El instrumento $Z$, al no estar correlacionado con $u$, nos permitir\'a utilizar solamente la parte ex\'ogena de $X$ para identificar su efecto en $Y$.
    \end{itemize}
\end{frame}

% \begin{frame}{Ejemplo: efecto del servicio militar sobre el crimen}
%     $$
%     \textit{crimen}=\alpha+\beta_{1} X+U
%     $$
%     $X$ toma valor $1$ si la persona realiz\'o el servicio militar en el pasado.
%     \begin{itemize}
%         \item $X$ es end\'ogena, $E[U \mid X] \neq 0$. \textbf{\textquestiondown Por qu\'e}?
%         \item Por lo tanto, el estimador MCO est\'a sesgado.
%         \item Instrumento $Z$: indicador igual a $1$ si la persona sali\'o sorteada aleatoriamente para hacer el servicio militar. \textbf{\textquestiondown Cumple los supuestos}?
%             \begin{enumerate}
%                 \item \textbf{Relevancia:} \underline{S\'\i}. $\operatorname{Cov}(Z, X) \neq 0$ porque, si eras sorteado, el Estado te obligaba a ir.
%                 \item \textbf{Exogeneidad:} \underline{S\'\i}. $\operatorname{Cov}(Z, U)=0$ porque $Z$ depende de un sorteo. No se correlaciona con nada m\'as que con la probabilidad de realizar el servicio militar.
%             \end{enumerate}
%         \item Entonces, podemos usar $Z$ para identificar el efecto. El m\'etodo utilizar\'a la variabilidad de las personas que realizaron el servicio militar \emph{solo} porque fueron sorteadas, y no la de quienes lo eligieron voluntariamente (end\'ogena).
%     \end{itemize}
% \end{frame}

\begin{frame}{Buenos y malos instrumentos}
    \begin{itemize}
        \item \textbf{\textquestiondown Cu\'al de los siguientes $Z$ es un buen instrumento}?
        \begin{figure}
            \centering
            \includegraphics[width=0.8\linewidth]{complemetario/Imagenes/Grafico4_3.jpeg}
        \end{figure}
        \item S\'olo el gr\'afico A muestra un instrumento v\'alido, ya que cumple con: \\
        1) \textbf{Relevancia:} $\operatorname{Cov}(z,x) \neq 0$ y 2) \textbf{Exogeneidad:} $\operatorname{Cov}(z,u)=0$.
        \item El caso B viola la exogeneidad porque el instrumento se correlaciona con el error ($\operatorname{Cov}(z,u) \neq 0$). 
        \item El caso C viola la relevancia. Si bien el instrumento no se correlaciona con el error, tampoco lo hace con $X$ ($\operatorname{Cov}(x,z)=0$). 
    \end{itemize}
\end{frame}

\begin{frame}{Ejemplos de variables end\'ogenas e instrumentos en la literatura}
    \begin{itemize}
        \item \textbf{Retornos de la educaci\'on:}
        $$
        \log \bigl(\text{salario}\bigr)_{i}=\beta_{0}+\beta_{1} \,\text{educ}_{i}+u_{i}.
        $$
        \begin{itemize}
            \item $Z$: trimestre de nacimiento (Angrist y Krueger, 1991).
            \item $Z$: proximidad geogr\'afica a una universidad (Card, 1995).
        \end{itemize}

        \item \textbf{Efecto de las instituciones en el PIB per c\'apita:}
        $$
        \text{GDP}_{2015,i}=\beta_{0}+\beta_{1}\, \text{instituciones}_{2015,i}+u_{i}.
        $$
        \begin{itemize}
            \item $Z$: mortalidad por enfermedades en el pa\'\i s cuando llegaron los europeos siglos atr\'as (Acemoglu \emph{et al.}, 2001).
        \end{itemize}
    \end{itemize}
\end{frame}

\section{Variables instrumentales}

\begin{frame}{Implementaci\'on del m\'etodo IV}
    \begin{itemize}
        \item M\'etodos para estimar $\hat{\beta}^{IV}$:
        \begin{enumerate}
            \item Estimaci\'on por Variables Instrumentales (VI) \textit{per se}.
            \item M\'\'inimos Cuadrados en Dos Etapas (MC2E), 
        \end{enumerate}
    \end{itemize}
\end{frame}

\begin{frame}{Estimaci\'on por Variables Instrumentales (VI)}
    \begin{itemize}
        \item Consideremos el modelo de regresi\'on simple:
        $$
        y_{i}=\beta_{0}+\beta_{1} x_{i}+u_{i}, \quad i=1,\ldots,N.
        $$
        \item Problema: $\operatorname{Cov}\left(u_{i},x_{i}\right) \neq 0$.
        \item Contamos con un instrumento $z$ que cumple:
        \begin{enumerate}
            \item $\operatorname{Cov}\left(u_{i},z_{i}\right)=0$.
            \item $\operatorname{Cov}\left(x_{i},z_{i}\right)\neq 0$.
        \end{enumerate}
        \item Entonces:
        $$
        \begin{aligned}
        \operatorname{Cov}\left(y_{i},z_{i}\right) &= \operatorname{Cov}\bigl(\underbrace{\beta_{0}+\beta_{1}x_{i}+u_{i}}_{y_{i}}, z_{i}\bigr) \\
        &= \underbrace{\operatorname{Cov}\left(\beta_{0},z_{i}\right)}_{0} + \beta_{1}\, \operatorname{Cov}\left(x_{i},z_{i}\right) + \underbrace{\operatorname{Cov}\left(u_{i},z_{i}\right)}_{0} \\
        &= \beta_{1}\, \operatorname{Cov}\left(x_{i},z_{i}\right).
        \end{aligned}
        $$
        \item Despejando:
$\beta_{1}=\frac{\operatorname{Cov}\left(y_{i},z_{i}\right)}{\operatorname{Cov}\left(x_{i},z_{i}\right)}.
        $
       
    \end{itemize}
\end{frame}

\begin{frame}{Relaci\'on poblacional y an\'alogo muestral}
    \begin{itemize}
     \item La {\color{blue}identificación} del
 parámetro en este contexto significa que se puede escribir $\beta_1$ en términos de los momentos
poblacionales que se pueden estimar mediante una muestra de datos.
      %  \item Con la relaci\'on a nivel \textbf{poblacional} podemos aplicar un criterio de analog\'\i a en la muestra y obtener un \textbf{estimador} de $\beta$:
        $$
        \hat{\beta}_{VI}=\frac{S_{y_{i},z_{i}}}{S_{x_{i},z_{i}}}=\frac{\sum_{i=1}^{N}(y_{i}-\bar{y})(z_{i}-\bar{z})}{\sum_{i=1}^{N}(x_{i}-\bar{x})(z_{i}-\bar{z})}.
        $$
        \item Todo esto es observable, por lo que podemos estimarlo (\texttt{R} lo hace por ustedes).
        \item Adem\'as, si se cumplen los supuestos sobre el instrumento $Z$, el estimador es consistente:
        $$
            \hat{\beta}_{VI} \xrightarrow{p} \beta.
        $$
    \end{itemize}
\end{frame}

\begin{frame}{Estimaci\'on con VI en modelos de regresi\'on m\'ultiple}
    \begin{itemize}
        \item La estimaci\'on por IV se extiende f\'acilmente a modelos con m\'\'ultiples regresores.
        \item Nos centraremos en el caso con un \textbf{\'unico} regresor end\'ogeno y varios controles ex\'ogenos, aunque es posible tener m\'as de una variable end\'ogena.
    \end{itemize}

    \underline{\textbf{Estimaci\'on por VI con una end\'ogena y varias ex\'ogenas}}\\[4pt]
    $$
    \ln \bigl(\text{salario}_{i}\bigr)=\beta_{0}+\beta_{1}\, \text{educ}_{i}+\beta_{2}\, \text{mujer}_{i}+u_{i}.
    $$
    \begin{itemize}
        \item Idea: usar las ex\'ogenas como sus propios instrumentos; la matriz $Z$ incluir\'a (i) la constante; (ii) el instrumento para \textit{educ}; y (iii) \textit{mujer}.
        \item Aplicamos la f\'ormula anterior:
        $$
        \hat{\beta}_{VI}=(Z^{T}X)^{-1}Z^{T}Y,
        $$
        donde $Z$ es de $N\times k+1$, $X$ es de $N\times k+1$, $Y$ es de $N\times 1$ y $\hat{\beta}^{IV}$ de $k+1\times 1$.
    \end{itemize}
\end{frame}

% \begin{frame}{Comparaci\'on entre $\hat{\beta}^{IV}$ y $\hat{\beta}^{MCO}$}
%     $$
%     \ln \bigl(\text{salario}_{i}\bigr)=\beta_{0}+\beta_{1}\, \text{educ}_{i}+u_{i}.
%     $$
%     \begin{itemize}
%         \item Supongamos que la \'unica variable omitida es la inteligencia.
%         \item Como la inteligencia se correlaciona positivamente con la educaci\'on y con el salario, esperar\'\i amos:
%         $$
%         E\left[\hat{\beta}^{MCO}\right]>\beta.
%         $$
%         \item Si disponemos de un instrumento v\'alido, entonces:
%         $$
%         E\left[\hat{\beta}^{IV}\right]=\beta.
%         $$
%         \item De este modo, al realizar ambas estimaciones esperar\'\i amos observar:
%         $$
%         \hat{\beta}^{MCO}>\hat{\beta}^{IV},
%         $$
%         porque corregimos el sesgo de MCO que, en este caso, sobreestimaba el efecto.
%         \item \textbf{\textquestiondown Qu\'e ocurre si no hay variable omitida, pero \textit{educ} est\'a mal medida}? En ese caso $\hat{\beta}^{MCO}<\hat{\beta}^{IV}$, porque eliminamos el sesgo de atenuaci\'on propio de MCO.
%     \end{itemize}
% \end{frame}

\begin{frame}{Resumen}
    \begin{itemize}
        \item Decimos que una variable es end\'ogena cuando $\operatorname{Cov}(u,x) \neq 0$.
        \item En presencia de endogeneidad se viola uno de los supuestos de Gauss--Markov, y $\hat{\beta}_{\text{MCO}}$ ser\'a:
        \begin{itemize}
            \item Sesgado.
            \item Inconsistente.
            \item No eficiente.
        \end{itemize}
        \item Existen tres fuentes habituales de endogeneidad en econom\'\i a:
        \begin{enumerate}
            \item Omisi\'on de variables relevantes.
            \item Simultaneidad.
            \item Errores de medida.
        \end{enumerate}
        \item Para estimar de manera consistente necesitaremos un \textbf{instrumento}.
        \item Un instrumento v\'alido es una tercera variable $z$ que cumple:
        \begin{itemize}
            \item \textbf{Exogeneidad:} $\operatorname{Cov}(u,z)=0$.
            \item \textbf{Relevancia:} $\operatorname{Cov}(x,z) \neq 0$.
        \end{itemize}
        \item El m\'etodo de IV permite estimar $\beta$ de forma consistente.
    \end{itemize}
\end{frame}


\begin{frame}{Inferencia y estimación por Variables Instrumentales}
\begin{itemize}
    \item Al igual que MCO, VI es un método de estimación
    \begin{itemize}
        \newline → Un método para obtener un estimador de nuestro parámetro de interés $\beta$
    \end{itemize}
    \item Tal como con MCO, una vez obtenidas las estimaciones \(\hat{\beta}_{VI}\) queremos hacer inferencia estadística sobre las mismas. Necesitamos errores estándar.
    \begin{itemize}
        \newline →  Ej. pruebas t de significación de los parámetros
    \end{itemize}
    \item Tal como en MCO, vamos a asumir homoscedasticidad, ahora condicional en el instrumento \(z\): \( \text{Var}(u \mid z) = \sigma^2 \)
    % \item Se puede demostrar que se puede estimar la varianza de \(\hat{\beta}_{VI}\) como:
    % \[
    % \text{Var}(\hat{\beta}_{VI}) = \frac{\hat{\sigma}^2}{\text{SST}_x R^2_{x,z}}
    % \]
\end{itemize}
\end{frame}
\begin{frame}{Inferencia y estimación por Variables Instrumentales}
\begin{itemize}
    \item Se puede demostrar que se puede estimar la varianza de \(\hat{\beta}_{VI}\) como:
    \[
    \text{Var}(\hat{\beta}_{VI}) = \frac{\hat{\sigma}^2}{\text{SST}_x R^2_{x,z}}
    \]
    donde
    \begin{enumerate}
        \item \(\hat{\sigma}^2 = \frac{1}{n-2} \sum_{i=1}^{n} \hat{u}_i^2\)
        \item \(R^2_{x,z}\) es el \(R^2\) de la regresión: \(X_i = \alpha_0 + \alpha_1 Z_i + v\)
        \item \(\text{SST}_x = \sum_{i=1}^{n} (x_i - \bar{x})^2\) es la suma de cuadrados total de \(X\).
    \end{enumerate}
\end{itemize}
\end{frame}
\begin{frame}{Estimador para la varianza de las estimaciones por VI}
\begin{itemize}
    \item Noten que:
    \[
    \text{Var}(\hat{\beta}_{VI}) = \frac{\hat{\sigma}^2}{\text{SST}_x R^2_{x,z}}
    \]
    \item Noto: Mientras más débil la relación entre \(x\) y su instrumento \(z\), menor \(R^2_{x,z}\)→ mayor \(\text{Var}(\hat{\beta}_{VI})\).
    \item l Modelo de Regresión Simple:
    \[
    \text{Var}(\hat{\beta}_{MCO}) = \frac{\hat{\sigma}^2}{\text{SST}_x}
    \]
    \item Dado que \(R^2_{x,z} \leq 1\), la varianza de las estimaciones por VI será mayor o igual a la varianza de los estimadores MCO.
    \begin{itemize}
        \item Mientras más fuerte la relación entre \(x\) y \(z\), más se acercarán las varianzas.
        \item Caso extremo: \(x\) es exógena, \(x\) es el instrumento de \(x\) → \(R^2_{x,z} = 1\) → \(\text{Var}(\hat{\beta}_{VI}) = \text{Var}(\hat{\beta}_{MCO})\).
    \end{itemize}
\end{itemize}
\end{frame}
\begin{frame}{Estimación por mínimos cuadrados en dos etapas}
\begin{itemize}
    \item Otro método para estimar consistentemente los parámetros con instrumentos es el Mínimos Cuadrados en 2 etapas (MC2E). El coeficiente que obtenemos será el mismo.
    \item Comenzamos con el siguiente modelo donde estamos interesados en el efecto de \( x_1 \) en \( y_1 \):
    \[
    y_i = \beta_0 + \beta_1 x_{1i} \text{(endo)} + \beta_2 x_{2i} \text{(exog)} + u_i \quad \forall i = 1, \ldots, N
    \]
    Con: 
\begin{enumerate}[(i)]
    \item \( \text{Cov}(x_{1i}, u_i) \neq 0 \),
    \item \( \text{Cov}(x_{2i}, u_i) = 0 \),
    \item \( \text{Cov}(z_i, u_i) = 0 \)
\end{enumerate}
    \item Por ejemplo:
    \[
    \ln(\text{salario}_i) = \beta_0 + \beta_1 \text{educ}_i + \beta_2 \text{mujer} + u_i
    \]
    \item Y podemos usar como instrumento de \textit{educ} la distancia al establecimiento educativo más cercano.
\end{itemize}
\end{frame}
\begin{frame}{MC2E - Etapa 1}

Primera etapa: $Z \rightarrow X$ Regresamos la variable endógena sobre todas las exógenas incluido el instrumento

\begin{itemize}
    \item \textbf{Paso 1:} Estimamos la siguiente ecuación por MCO.
    \begin{equation}
    x_{1i} = \beta_0 + \beta_1 z_i + \beta_2 x_{2i} + v_i
    \end{equation}

    \item \textbf{Paso 2:} Obtenemos los $\hat{\beta}$ y construimos el valor predicho de $x_1$:
    \begin{equation}
    \hat{x}_{1i} = \hat{\beta}_0 + \hat{\beta}_1 z_i + \hat{\beta}_2 x_{2i}
    \end{equation}

    \item Noten que $\hat{x}_1$ es una función de $z$ y $x_2$ que son variables exógenas. Por lo tanto, no correlaciona con el error $(u)$. El instrumento permite quedarnos con la parte exógena de $x$.

    \item \textbf{Intuición:} $x_i = \hat{x}_i + \hat{v}_i$. Solo usamos la parte exógena: $\hat{x}_i$ y nos deshacemos de la endógena.

    \item \textbf{Paso 3:} Noten que durante este paso, se aprovecha para testear la relevancia del instrumento. Haciendo un test de $\hat{\beta}_1 \neq 0$. 
\end{itemize}

\end{frame}
\begin{frame}{MC2E - Etapa 2}

Segunda etapa: $\hat{x} \rightarrow y$

\begin{itemize}
    \item En la ecuación original, reemplazamos la variable endógena $x_1$ por la predicción de esta variable que obtuvimos en la primera etapa $\hat{x}_1$ (Ec.2).
    
    \item Nuestra ecuación original (Ec.1) la podemos escribir entonces:
    \begin{equation}
    y_i = \beta_0 + \beta_1 \hat{x}_{1i} + \beta_2 x_{2i} + u_i
    \end{equation}
    
    \item $\hat{x}_{1i}$ es exógena, ya que es una combinación lineal de variables exógenas (Ec.2).
    
    \item Ahora, ¡podemos estimar Ec.3 por MCO!
    \begin{equation}
    \hat{\beta}_{MC2E} = (\hat{X}^T \hat{X})^{-1} \hat{X}^T Y
    \end{equation}
    
    \item Con $\hat{X} =    \begin{bmatrix}
    1 & \hat{x}_{1,1} & x_{2,1} \\
    1 & \hat{x}_{1,2} & x_{2,2} \\
    \vdots & \vdots & \vdots \\
    1 & \hat{x}_{1,n} & x_{2,n} \\
    \end{bmatrix}$ 
    
    \item El $\hat{\beta}_1$ hallado de esta manera se conoce como $\hat{\beta}_{1,MC2E}$.
\end{itemize}

\end{frame}
% \begin{frame}{MC2E - Varianza e Inferencia}

% \begin{itemize}
%     \item Si bien $\hat{\beta}_{MC2E}$ son la estimación por MCO de la segunda etapa, su varianza no es igual a la varianza de un estimador MCO ya que debe recoger las correlaciones entre la primera y la segunda etapa.

%     \item La varianza de los estimadores MC2E se puede aproximar asintóticamente a:
%     \begin{equation}
%     \text{Var}(\hat{\beta}_{1,MC2E}) \approx \frac{\sigma^2}{SST_{\hat{x}}(1 - R^2_1)} 
%     \end{equation}
% %\quad \text{Var}(\hat{\beta}_{1,MCO}) = \frac{\sigma^2}{SCT_{x}(1 - R^2_1)}
%     \item Donde:
%     \begin{itemize}
%         \item $SST_{\hat{x}} = \sum_{i=1}^{n} (\hat{x}_i - \bar{\hat{x}})^2$ es la variación total de $\hat{x}$.
%         \item $R^2_1$ es el $R^2$ de regresar $\hat{x}_1$ sobre los demás regresores $(x_2)$.% e.g., obtener el $R^2$ de $\hat{x}_1 = \beta + \beta_1 x_2 + c$.
%     \end{itemize}

%     \item Siempre $SST_{x} \geq SST_{\hat{x}}$. ¿Intuición? Recuerden que de toda la variabilidad de $X$, agarramos solo una parte (la parte exógena). %Además, generalmente, $R^2_1 \geq R^2$.
    
%     \item Entonces: $\text{Var}(\hat{\beta}_{1,MC2E}) \geq \text{Var}(\hat{\beta}_{1,MCO})$.

%     \item Al igual que con VI, una relación débil entre la variable endógena y su instrumento incrementa la varianza de la estimación \textcolor{red}
% {(porque achica el $SST_{\hat{x}}$).}
% \end{itemize}

% \end{frame}

% \begin{frame}{MC2E - Ejemplo 1 Etapa}

% \begin{itemize}
%     \item Queremos saber los efectos del turismo en México en los ingresos de la ciudad.
%     \begin{equation}
%     \ln(\text{income}_i) = \beta_0 + \beta_1 \text{HotelSales}_i + \beta_2 \log(\text{temp}_i) + u_i
%     \end{equation}

%     \item Ciudades más ricas pueden atraer más turismo, entonces hay endogeneidad.

%     \newline Usamos: distancia de la ciudad a USA como instrumento y estimamos por MC2E.
% \end{itemize}

% \end{frame}
% \begin{frame}{Primera Etapa}

% \begin{itemize}
%     \item Estimamos la primera etapa por MCO (instrumentos en $X$)
%     \begin{equation}
%     \log(\text{HotelSales}_i) = \beta_0 + \beta_1 \text{distUS}_i + \beta_2 \log(\text{temp}_i) + v_i
%     \end{equation}

%                      \begin{figure}
%         \centering
%         \includegraphics[width=0.8\linewidth]{clase18/images/1.png}
%     \end{figure}
    
% \end{itemize}

% \end{frame}
% \begin{frame}{MC2E - Ejemplo 2 Etapa: Predicción de $X$ en $Y$}

% \begin{itemize}
%     \item Usamos los valores predichos $\widehat{\text{HotelSales}}$ con la primera etapa y estimamos la segunda etapa:
%     \begin{equation}
%     \ln(\text{Income}_i) = \beta_0 + \beta_1 \log(\widehat{\text{HotelSales}}_i) + \beta_2 \log(\text{temp}_i) + u_i
%     \end{equation}

%                      \begin{figure}
%         \centering
%         \includegraphics[width=0.8\linewidth]{clase18/images/2.png}
%     \end{figure}
    
% \end{itemize}

% \end{frame}
% \begin{frame}{MC2E - Comparamos con MCO}

% \begin{itemize}
%     \item 0.45 con una estimamos por MCO:

%                  \begin{figure}
%         \centering
%         \includegraphics[width=0.8\linewidth]{clase18/images/3.png}
%     \end{figure}
    
%     \item Estimación por MC2E: un incremento de 1\% en las ventas de hoteles aumenta 0.45\% el ingreso de la ciudad.
%     \item Estimación por MCO: un incremento de 1\% de las ventas en hoteles aumenta 0.05\% el ingreso en la ciudad. MCO estaba sesgado, subestimando el verdadero valor.
% \end{itemize}

%\end{frame}
\begin{frame}{Conclusiones}

\begin{itemize}
    \item Vimos dos métodos para estimar consistentemente los parámetros cuando el regresor es endógeno y contamos con un instrumento válido:
    \begin{itemize}
        \item VI
        \item MC2E
    \end{itemize}
    
    \item En el Modelo de Regresión Simple (MRS), con un solo instrumento, ambos métodos nos darán las mismas estimaciones puntuales y los mismos errores estándar.
    
    \item Los resultados serán distintos cuando tengamos más de un instrumento para la variable endógena.
    
    \item Estimar por estos métodos nos asegura la consistencia   de las estimaciones siempre que los instrumentos sean válidos.
    
    \item Pero pagaremos el precio de menos precisión en esas estimaciones (errores estándar más grandes que en MCO).
    
   
\end{itemize}

\end{frame}
\end{document}
